\index{Repère cartésien}\index{Coordonnées cartésiennes}
Pour repérer la position du \robot\ sur l'écran, nous utilisons un {\em repère cartésien}
tel qu'il est représenté à la figure \ref{fig:repcar} ci-dessous~:

\begin{figure}[htbp]
 \includevspacePostScript[6.2cm]{RepereCartesien}{0.95}{0.95}
 \caption{Coordonnées dans un repère cartésien}\label{fig:repcar}
\end{figure}

L'{\em origine\/} $0$ du repère est située au centre de la grille. L'axe horizontal $OX$ est 
appelé l'axe des {\em abscisses\/}\index{Abscisse}. L'axe vertical $OY$ est 
appelé l'axe des {\em ordonnées\/}\index{Ordonnée}. Ces deux axes sont gradués. La longueur des
graduations sur l'axe des abscisses est égale à la longueur du côté horizontal d'un carreau 
 de la grille comptée en pixels. De même, la longueur des
graduations sur l'axe des ordonnées est égale à la longueur du côté vertical d'un carreau
de la grille comptée en pixels. 

Un point de la grille est repéré par ses {\em  coordonnées cartésiennes}\index{Coordonnées cartésiennes} qui sont
deux entiers $(X, Y)$, le premier étant son {\em abscisse} et le second son {\em ordonnée}. Par exemple le 
point de la grille situé au centre de la grille a pour coordonnées $(0, 0)$. Les points situés
à droite de l'origine $O$ ont une abscisse positive. Les points situés
à gauche de l'origine $O$ ont une abscisse négative (précédée du signe moins $-$). 
De même, les points situés au dessus
de l'origine $O$ ont une ordonnée positive tandis que les points situés
en dessous de l'origine $O$ ont une ordonnée négative.  A titre d'exemple, quelques points de 
la grille sont repérés sur la figure \ref{fig:repcar}  (page \pageref{fig:repcar}) par leurs 
coordonnées.  

L'abscisse X d'un point A est la distance à l'origine de sa projection sur l'axe $OX$
(c'est-à-dire du point d'intersection
de l'axe $OX$ avec la droite parallèle à l'axe $OY$ passant par A). De même, l'ordonnée Y d'un point 
A est la distance à l'origine de sa projection sur l'axe $OY$
(c'est-à-dire du point d'intersection
de l'axe $OY$ avec la droite parallèle à l'axe $OX$ passant par A). Par exemple sur la figure
\ref{fig:coorA}, 
ci-dessous,
on a $X = 6$ et $Y = 7$\footnote{\ On trouvera sur la disquette d'accompagnement du livre un programme
\pascal\ permettant d'apprendre à repérer des points du plan avec
un repère cartésien.}.

\begin{figure}[htbp]
 \includevspacePostScript[4.1cm]{AbscisseOrdonnee}{0.95}{0.95}
 \caption{Point A d'abscisse X = 6 et d'ordonnée Y = 7}
 \label{fig:coorA}
\end{figure}

La commande {\idt deplacer(X, Y)} place le \robot\ sur le point de la grille de coordonnées cartésiennes $(X, Y)$. 
L'orientation et la position du crayon du \robot\ restent inchangées. 

Par exemple, la commande {\idt deplacer(0, 0)} place le \robot\ à l'origine du repère
cartésien au centre de la grille\footnote{\ La commande {\idt ce} a l'effet supplémentaire
de recentrer si nécessaire l'origine du repère exactement au centre de l'écran.}. 
Un exemple est donné ci-après~:

\clearpage
\vspace*{-1cm}
 {\raggedright 
 \vbox to 0pt{\vspace*{\baselineskip}\vspace*{6cm}%
 \includevspacePostScript[-\baselineskip]{LigneDeCarres}{0.95}{0.95}}%
 \label{prog:LigneDeCarres1}
 \input LigneDeCarres1.tex
 }

La commande {\idt translater(X, Y)} permet de translater le \robot\ de $(X, Y)$
sans changer son orientation\footnote{\ On trouvera sur la disquette d'accompagnement du livre un programme
\pascal\ permettant d'apprendre à translater des points du plan 
repérés par des coordonnées cartésiennes.}. Un exemple de translation d'un rectangle est donné à la figure 
\ref{fig:extrans} ci-dessous.
\begin{figure}[h]
\includevspacePostScript[6.0cm]{Translater}{0.95}{0.95}%
 \caption{Translation (17, -8) de A en B}\label{fig:extrans}
\end{figure}
\newpage
Une translation du \robot\ de $(X, Y)$ consiste à augmenter (ou diminuer si X est négatif)
son abscisse de X et à augmenter (ou diminuer si Y est négatif)
son ordonnée de Y. Si le centre de la grille de coordonnées $(0, 0)$ est exactement au centre
de l'écran, on peut définir le déplacement à l'aide de la translation comme suit~:

{
\parindent \parindentprog
\medskip

\vspace{\neglskip}\mbox{\blc  \blk\blk\idt  \mcl procedure\idt \ deplacer$($X\ $:$\ integer$;$\ Y\ $:$\ integer$)$$;$}\par 
\vspace{\neglskip}\mbox{\blc  \blk\blk\idt  \mcl begin\idt }\par 
\vspace{\neglskip}\mbox{\blc  \blk\blk\blk\idt  ce$;$}\par 
\vspace{\neglskip}\mbox{\blc  \blk\blk\blk\idt  translater$($X$,$\ Y$)$$;$}\par 
\vspace{\neglskip}\mbox{\blc  \blk\blk\idt  \mcl end\idt $;$}\par 

\medskip }

Le programme {\idt LigneDeCarres1} de la page
\pageref{prog:LigneDeCarres1} peut être réécrit en uti\-li\-sant une
commande de translation au lieu d'une commande de dé\-pla\-ce\-ment
com\-me suit~:\\
{\raggedright%
\vspace{\baselineskip}%
\vbox to 0pt{\vspace*{\baselineskip}\vspace*{6cm}%
\includevspacePostScript[-\baselineskip]{LigneDeCarres}{0.95}{0.95}}%
\vspace*{-1cm}%
\input LigneDeCarres2.tex
}
 
 Avec la commande {\idt translater(X, Y)}, le déplacement du \robot\ est relatif à sa position de
départ tandis qu'avec la commande {\idt deplacer(X, Y)}, le déplacement du \robot\ est relatif 
à l'origine du repère cartésien.
 
\noindent \rule{137mm}{0.1mm} %---------------------------------------------
\begin{exercices}{Écrire des programmes \pascal\  utilisant les commandes
{\idt de\-pla\-cer} et {\idt translater} pour réaliser les dessins suivants~:

\tiret  les pavages obtenus par répétition de divers motifs inscrits dans un carré de
dimension $(6 \times 6)$ de la figure \ref{fig:pm6x6} (page \pageref{fig:pm6x6})$\:$;

\tiret  le pavage de quadrilatères de la figure \ref{fig:pavqua} (page \pageref{fig:pavqua})$\:$;

\tiret  la mosaïque de l'Alhambra relevée par Escher \cite{escher1} de la figure 
\ref{fig:mosa5} (page \pageref{fig:mosa5}).

\noindent $\Box$}\end{exercices}
\begin{figure}[htbp]
 \includevspacePostScript[7.7cm]{PavageQuadrilateres}{0.95}{0.95}%
 \caption{Pavage de quadrilatères}\label{fig:pavqua}
\end{figure}
\begin{figure}[htbp]
 \includevspacePostScript[5.8cm]{MosaiqueAlhambra5}{0.95}{0.95}
 \caption{Mosaïque de l'Alhambra}\label{fig:mosa5}
\end{figure}
\begin{figure}[htbp]
 \includevspacePostScript[17.2cm]{PavageMotifs6x6}{0.95}{0.95}
 \caption{Pavages de motifs $(6 \times 6)$}\label{fig:pm6x6}
\end{figure}
\begin{corrige}{
Nous donnons le programme de dessin du pavage de quadrilatères de la figure \ref{fig:pavqua} (page 
\pageref{fig:pavqua}), les autres programmes se trouvant sur disquette. Il s'agit d'un pavage de
l'écran avec des motifs inscrits dans des carrés $(6 \times 6)$, le motif de base étant
dessiné à côté de la procédure qui le dessine.
%\noindent\includePostScript{MotifPavageQuadrilateres}{0.95}{0.95}
\input PavageQuadrilateres.tex
\noindent $\Box$}\end{corrige}

%%%%%%%%%%%%%%%%%%%%%%%%%%%%%%%%%%%%%%%%%%%%%%%%%%%%%%%%%%%%%%%%%%%%%%%%%%%
\noindent \framebox[\textwidth]{
\begin{minipage}{\myparboxwidth}
\setlength{\parindent}{2em}
\small \sf
\vspace{\abovedisplayskip}
\section*{Addenda}
 
 Suivant les traces d'Apollonius\footnote{\ \href{https://fr.wikipedia.org/wiki/Apollonios_de_Perga}{Apollonius de Perge}, astronome et
 mathématicien grec, 295 -- 230 av. J.-C.} qui avait donné informellement les équations
 des coniques (cercles, ellipses, paraboles et hyperboles), Fermat et Descartes ont  été les premiers 
 à ``mettre en équations'' des problèmes géométriques. 
 Descartes a introduit la notion
 de système de coordonnées planes, chaque point du plan étant défini par son abscisse $x$
 et son ordonnée $y$ obtenues en projetant sur des axes $OX$ et $OY$, rectangulaires ou obliques. 
 
 Une courbe algébrique est alors définie par la relation $f(x, y) = 0$ qui existe entre les coordonnées de
 chacun de ses points. 
Par exemple l'équation cartésienne du cercle de centre $(a, b)$ est
 $x^{2} + y^{2} - 2 a x - 2 b y + c = 0$ avec le rayon $R$ tel que $R^{2} = a^{2} + b^{2} -c$
 en supposant que $a^{2} + b^{2} \geq c$. 
 
\vspace{\abovedisplayskip}
\end{minipage}
}
%%%%%%%%%%%%%%%%%%%%%%%%%%%%%%%%%%%%%%%%%%%%%%%%%%%%%%%%%%%%%%%%%%%%%%%%%%%% 

%%%%%%%%%%%%%%%%%%%%%%%%%%%%%%%%%%%%%%%%%%%%%%%%%%%%%%%%%%%%%%%%%%%%%%%%%%%
\noindent \framebox[\textwidth]{
\begin{minipage}{\myparboxwidth}
\setlength{\parindent}{2em}
\small \sf
\vspace{\abovedisplayskip}
\indent  Descartes a exposé sa méthode dans un ouvrage
 publié en 1637 et
 intitulé ``La Géométrie'' qui, avec ``La Dioptrique'' et ``Les Météores'', constituaient les
 appendices d'un des textes les plus célèbres de la philosophie, le fameux ``Discours de la 
 méthode pour bien conduire sa raison et chercher la vérité 
 dans les sciences'' où il expose sa philosophie méthodique et rationelle appelée depuis le
{\em  cartésianisme}. L'idée d'équation apparaît plus clairement chez Fermat, qui a découvert
 indépendamment de Descartes, vers 1629, le principe de base de la {\em géométrie
 analytique}. Il ne l'a publié qu'en 1679 dans son ``Isagoge (Introduction aux lieux plans et solides)''.
 
 Descartes n'avait étudié sérieusement que les coniques du second degré. 
 Newton\footnote{\ \href{https://fr.wikipedia.org/wiki/Isaac_Newton}{Isaac Newton}, mathématicien, physicien et astronome anglais, 1642--1727.}
 étudia les cubiques du troisième degré et les courbes de degré supérieur. On lui
 doit aussi d'avoir employé systématiquement les coordonnées négatives dont Descartes se
 méfiait. La Hire\footnote{\ \href{https://fr.wikipedia.org/wiki/Philippe_de_La_Hire}{Philippe de la Hire}, astronome et mathématicien français, 1640--1718}
 introduisit en 1679 la considération d'un système de coordonnées spatiales, ce qui
 revient à construire un  troisième axe de coordonnées $OZ$, perpendiculaire aux axes $OX$
 et $OY$ définissant le plan. Un point dans l'espace est alors déterminé par trois nombres
 réels $(x, y, z)$ et l'équation $f(x, y, z) = 0$ définit une surface gauche. 
 Möbius\footnote{\ \href{https://fr.wikipedia.org/wiki/August_Ferdinand_Möbius}{August Ferdinand Möbius}, astronome et mathématicien allemand,
 1790--1868} et Plücker%
 \footnote{\ \href{https://fr.wikipedia.org/wiki/Julius_Plücker}{Julius Plücker}, mathématicien et physicien allemand, 1801--1868} étendirent
 la notion de coordonnées en introduisant les coordonnées homogènes $(X, Y, Z, T)$ telles que
 $x = \frac{X}{T}$,  $y = \frac{Y}{T}$ et  $z = \frac{Z}{T}$. Ceci permet d'introduire les points
 à l'infini $(X, Y, Z, 0)$ et les points imaginaires $X = A_{1} + i A_{2}$, $Y = B_{1} + i B_{2}$,
 $Z = C_{1} + i C_{2}$ et $T = D_{1} + i D_{2}$ où $i^{2} = -1$. On peut alors considérer par
 exemple des cercles imaginaires tels que $a^{2} + b^{2} < c$.
 
 \vspace{\abovedisplayskip}
\end{minipage}
}
%%%%%%%%%%%%%%%%%%%%%%%%%%%%%%%%%%%%%%%%%%%%%%%%%%%%%%%%%%%%%%%%%%%%%%%%%%%% 
\newpage                               % Si on termine le chapitre par une page blanche
\thispagestyle{empty}           % Si on termine le chapitre par une page blanche

